\subsection{Differential operators}

Let \(n \in N\) and U an open subset of \(R^{n}\) . We equip \(R^{n}\) and U with the Lebesgue measure denoted \(\mu\) .

Let \(k \in N\) and \(h \in N\) such that \(h \geqslant k\) . Let D be a scalar differential operator on U of order \(\leqslant k\) , with coefficients \((n_{\alpha})_{|\alpha| \leqslant k}\) of class \(C^{h}\) on U. Suppose that for every \(\alpha\) such that \(|\alpha| \leqslant k\) and every \(\beta\) such that \(0 \leqslant \beta \leqslant \alpha\) , the function \(\partial^{\beta} n_{\alpha}\) is bounded on U.

Let \({}^{t}D\) the scalar differential operator on U transpose of D (VAR, R2, 14.3.2); it is of order \(\leqslant k\) and of class \(C^{h-k}\) ; for \(\varphi\in\mathcal{D}(U)\) , we have

\[
{ } ^ { t } \mathrm{D} ( \varphi ) = \sum _ { | \alpha | \leqslant k } ( - 1 ) ^ { | \alpha | } \partial ^ { \alpha } ( \overline { { n } } _ { \alpha } \varphi )
\]

(loc. cit.); in particular, the coefficients of \({}^{t}D\) are bounded on U.

We denote by D\_ the partial operator on \(L^{2}(U)\) with domain \(\mathcal{D}(U)\) defined by

\[
\varphi \mapsto \mathrm{D} (\varphi) = \sum_ {| \alpha | \leqslant k} n _ {\alpha} \partial^ {\alpha} \varphi .\tag{5}
\]

Let \(H_{D}\) the space of \(f \in L^{2}(U)\) such that the distribution

\[
\mathrm{D} (f) = \sum_ {| \alpha | \leqslant k} n _ {\alpha} \partial^ {\alpha} f
\]

belongs to \(L^{2}(U)\) ; we denote by \(D_{+}\) the partial operator with domain \(H_{D}\) defined by \(f \mapsto \mathrm{D}(f)\) .

Since we have \(\partial^{\alpha}f\in \mathrm{L}^{2}(\mathrm{U})\) if \(f\in \mathrm{H}^k (\mathrm{U})\) and \(|\alpha |\leqslant k\) , the Sobolev space \(\mathrm{H}^k (\mathrm{U})\) is contained in \(\mathrm{H}_{\mathrm{D}}\) ; in general, these spaces are distinct.

One has \(D_{-} \subset D_{+}\) , and these are differential operators associated with D on \(L^{2}(U)\) (def. 5 of IV, p. 235).

PROPOSITION 13. --- Let u be a partial operator on L²(U). If D₋ ⊂ u ⊂ D₊, then u is closable and (ᵗD)₋ ⊂ u* ⊂ (ᵗD)₊.

Let \(\varphi\) and \(\psi\) in \(\mathcal{D}(\mathrm{U})\) . We then have

\[
\begin{array}{c} \langle \varphi | \mathrm{D} (\psi) \rangle = \sum_ {| \alpha | \leqslant k} \int_ {\mathrm{U}} n _ {\alpha} \overline {{\varphi}} \partial^ {\alpha} \psi d \mu \\ = \sum_ {| \alpha | \leqslant k} (- 1) ^ {| \alpha |} \langle \partial^ {\alpha} (\overline {{n}} _ {\alpha} \varphi) | \psi \rangle = \langle {} ^ {t} \mathrm{D} (\varphi) | \psi \rangle \end{array}
\]

(cf.VAR, R2, 14.3.8). Since \(\mathcal{D}(\mathrm{U})\) is dense in \(\mathrm{L}^2 (\mathrm{U})\) this implies that \(\varphi \in \mathrm{dom}(u^{*})\) and \(u^{*}(\varphi) = {}^{t}\mathrm{D}(\varphi)\) . We therefore have \(({}^{t}\mathrm{D})_{-}\subset u^{*}\) In particular, \(u^{*}\) is densely defined and \(u\) is closable (prop. 8 of IV, p. 237).

Let \(f \in \mathrm{dom}(u^{*})\) and \(\varphi \in \mathcal{D}(\mathrm{U})\) . Since \(\mathcal{D}(\mathrm{U}) \subset \mathrm{dom}(u)\) , the distribution associated with \(u^{*}(f)\) satisfies

\[
\langle u ^ {*} (f), \varphi \rangle = \langle \overline {{\varphi}} | u ^ {*} (f) \rangle = \langle u (\overline {{\varphi}}) | f \rangle .
\]

Since \(D_{-} \subset u\) , one computes \(u(\overline{\varphi})\) by the formula (5), whence

\[
\begin{array}{l} \langle u ^ {*} (f), \varphi \rangle = \sum_ {\alpha} \langle n _ {\alpha} \partial^ {\alpha} \overline {{\varphi}} | f \rangle = \sum_ {\alpha} \langle \partial^ {\alpha} \overline {{\varphi}} | \overline {{n}} _ {\alpha} f \rangle \\ \qquad = \sum_ {\alpha} \langle \overline {{n}} _ {\alpha} f, \partial^ {\alpha} \varphi \rangle = \sum_ {\alpha} (- 1) ^ {| \alpha |} \langle \partial^ {\alpha} (\overline {{n}} _ {\alpha} f), \varphi \rangle = \langle {} ^ {t} D (f), \varphi \rangle . \end{array}
\]

The distributions \(u^{*}(f)\) and \({}^{t}\mathrm{D}(f)\) are therefore equal; the distribution \(f\) therefore belongs to \(H_{tD}\) and \(u^{*}(f)={}^{t}\mathrm{D}(f)\) , whence \(u^{*}\subset({}^{t}\mathrm{D})_{+}\) .

Remark. --- The proposition means that the adjoint of a partial operator \(u\) such that \(\mathrm{D}_{-} \subset u \subset \mathrm{D}_{+}\) can always be computedd in the sense of distributions: the elements \(f\) of the domain of \(u^{*}\) are elements of \(\mathrm{L}^{2}(\mathrm{U})\) such that the distribution \(^{t}\mathrm{D}(f)\) belongs to \(\mathrm{L}^{2}(\mathrm{U})\) , and we have \(u^{*}(f) = ^{t}\mathrm{D}(f)\) .

We say that D is formally symmetric if \({}^{t}D = D\) as a scalar differential operator. If that is the case, the partial operator D\_ is symmetric.

Consider the particular case of the second-order scalar differential operator defined by

\[
\Delta = - \sum_ {i = 1} ^ {n} \partial_ {i} ^ {2}.
\]

A Laplacian on U is any partial operator u, self-adjoint on \(L^{2}(U)\) such that \(\Delta_{-} \subset u\) (cf.VAR, R2, 14.4.3, p. 83). We shall see later (IV, p. 261, example) that there always exists at least one Laplacian on \(L^{2}(U)\) ; there may exist more than one (exercise 17 of IV, p. 358).
% semantic-fragments: legacy-input-seam
\relax{}
